%==============================================================================
% Page Layout 
%==============================================================================
\ifdefined\pdfminorversion\pdfminorversion=7\fi % Allow inclusion of PDF figures up to v1.7 (pdfTeX only)
\usepackage[letterpaper]{geometry} % Letter paper size
\geometry{
    verbose, % Show verbose output
    tmargin=1in, % Top margin
    bmargin=1in, % Bottom margin 
    lmargin=1.2in, % Left margin
    rmargin=1.2in % Right margin
}

% Enable margin changes locally
\usepackage{changepage} % Allows the use of adjustwidth

\usepackage{microtype} % Try to fix justification
% Increase tolerance and emergency stretch to avoid off-margin words
\tolerance=500
\emergencystretch=2em

\setlength{\parindent}{10pt} 
\setlength{\parskip}{0pt}

\usepackage[utf8]{inputenc} % UTF-8 input encoding
\usepackage[T1]{fontenc}
\usepackage{lmodern} % improves appearance, esp. with [T1]{fontenc}
\usepackage{fancyhdr} % Custom headers and footers
\usepackage{setspace} % Custom line spacing

% Tweaking lists
\usepackage{enumitem}
% Enable \pageref{LastPage} to write page # of ##
\usepackage{lastpage}

\usepackage{lscape}
\usepackage{threeparttable}
%==============================================================================
% Tables
%==============================================================================
\usepackage{booktabs,tabularx} % Advanced tables

% ======================================================================
% Packages
% ======================================================================

% Import AMS math, mathtools, fonts, and symbols
\usepackage{amsmath, mathtools, amsfonts, amssymb, bbm}
\usepackage{amsthm}

  % Theorem environments used in the appendix (Proposition~\ref{prop:pooling})
  \theoremstyle{plain}
  \newtheorem{proposition}{Proposition}
  \newtheorem{lemma}{Lemma}
  % Custom style for assumptions: optional name in italics, no parentheses.
  % Renders as: "Assumption D.1 Observed-period orthogonality." (name italicized).
  \newtheoremstyle{namedassumption}{}{}{}{}{\bfseries}{.}{ }{\thmname{#1}\thmnumber{
  #2}\thmnote{ \normalfont\itshape #3}}
  \theoremstyle{namedassumption}
  \newtheorem{assumption}{Assumption}
  \theoremstyle{remark}
  \newtheorem{remark}{Remark}
% Enable UTF-8 input encoding
\usepackage[utf8]{inputenc}

% Enable graphics support and set graphics path
\usepackage{graphicx}
\graphicspath{ {figures/} }

% Enable URL support
\usepackage{url}

% Enable fancy headers
\usepackage{fancyhdr}

% Switch to pdflscape instead of lscape (informs pdf viewer that page is landscape)
\usepackage{pdflscape}

% Fancy looping for the appendix
\usepackage{pgffor}

% Enable advanced table features
\usepackage{booktabs, multicol, multirow, floatrow, ragged2e, rotating, adjustbox}
\floatsetup[table]{capposition=top}
\floatsetup[figure]{capposition=top}
\usepackage{float}
\setlength{\textfloatsep}{10pt}
\usepackage{caption}
\newcommand\fnote[1]{\floatfoot{\small #1}}
\usepackage{subcaption}
\raggedbottom
% Enable float barrier
\usepackage{placeins}
% Enable line spacing
\usepackage{setspace}

% Enable color support and define custom colors
\usepackage[dvipsnames]{xcolor}
\definecolor{DarkBlue}{HTML}{090994}

% Enable hyperlinks
\usepackage[plainpages=false,pdfpagelabels=true]{hyperref}
\hypersetup{
    colorlinks=true,
    allcolors=DarkBlue,
    allbordercolors=white,
    pdfpagemode=FullScreen,
    breaklinks=true,
    pdfhighlight=/N
}

%==============================================================================
% Bibliography
%==============================================================================
\usepackage[longnamesfirst,authoryear]{natbib} % Bibliography style
\setcitestyle{round} % Round parentheses in citations
\bibliographystyle{ecca}
% \bibliographystyle{ecta_YEAR}

% Change some distances in bibliography:
    \setlength{\bibsep}{3pt}
    \setlength{\bibhang}{7pt}

% ======================================================================
% Commands and Specifications
% ======================================================================

% Set paragraph indentation
\setlength{\parindent}{4ex}

% Custom footnotes
\newcommand{\footremember}[2]{%
   \footnote{#2}
    \newcounter{#1}
    \setcounter{#1}{\value{footnote}}%
}
\newcommand{\footrecall}[1]{%
    \footnotemark[\value{#1}]%
} 
% ======================================================================
% Change font size of abstract 
% ======================================================================
\renewenvironment{abstract}
 {\small
  \begin{center}
  \bfseries \abstractname \vspace{-.5em}\vspace{0pt}
  \end{center}
  \list{}{
    \setlength{\leftmargin}{.3in}  % Adjust this to make the abstract wider or narrower
    \setlength{\rightmargin}{\leftmargin}  % Keep this the same as leftmargin for symmetry
  }%
  \item\relax}
 {\endlist}


% ======================================================================
% Blockquote
% ======================================================================
\definecolor{blockquote-border}{RGB}{150,150,150}
\definecolor{blockquote-text}{RGB}{60,60,60}
\usepackage{mdframed}
    \newmdenv[rightline=false, 
            bottomline=false, 
            topline=false, 
            linewidth=3.5pt, 
            leftmargin=1em,
            linecolor=blockquote-border,
            skipabove=\parskip]{customblockquote}

\renewenvironment{quote}{\begin{customblockquote}\list{}{\rightmargin=0em\leftmargin=0em}%
\item\relax\color{blockquote-text}\ignorespaces}{\unskip\unskip\endlist\end{customblockquote}}

% Fix Stata esttab output \sym command
\newcommand{\sym}[1]{\rlap{#1}}

% =====================================================================
% GRC table macros (added 2026-04-28 for Phase 1b of pipeline refactor).
% Captions and notes for GRC tables now live here, not in 0_programs.do.
% Cell catalog and design notes:
%   docs/specs/2026-04-24_grc-pipeline-refactor.md (M3, M11)
%   quality_reports/reviews/2026-04-28_paper-table-text-extraction.md
% =====================================================================

\usepackage{etoolbox}  % provides \ifstrequal used by macros below

% --- Reused phrases ----------------------------------------------------
\newcommand{\seclusterind}{We report robust standard errors, clustered at the individual level, in parentheses.}
\newcommand{\sigstars}{Stars denote: $^{*} p<0.10$; $^{**} p<0.05$; $^{***} p<0.01$.}
\newcommand{\refdatasec}{Please refer to Section~\ref{sec:data} for further details on the data}
\newcommand{\refIDNcanon}{and to the notes of Table~\ref{tab:GRC_IDN_consumption_urban_unb} for additional information on the variables}

% --- Country data sources (TZA standardized on "Tanzanian National Panel Survey") ---
\newcommand{\datasourceIDNunb}{This table uses data from the Indonesia Family Life Survey.}
\newcommand{\datasourceIDNbal}{This table uses the balanced panel from the Indonesia Family Life Survey.}
\newcommand{\datasourceCHNunb}{This table uses data from the China Family Panel Survey.}
\newcommand{\datasourceCHNbal}{This table uses the balanced panel from the China Family Panel Survey.}
\newcommand{\datasourceTZAunb}{This table uses data from the Tanzanian National Panel Survey.}
\newcommand{\datasourceTZAbal}{This table uses the balanced panel from the Tanzanian National Panel Survey.}

% --- Caption components: defined with name suffixes that match macro args
% (lowercase arg passed to \csname directly, no case-translation needed)
\newcommand{\countrynameIDN}{Indonesia}
\newcommand{\countrynameCHN}{China}
\newcommand{\countrynameTZA}{Tanzania}
\newcommand{\treatmenturban}{Urban Location}
\newcommand{\treatmentnonag}{Non-Agricultural Sector}
\newcommand{\depvarconsumption}{Consumption}
\newcommand{\depvarincome}{Income}
\newcommand{\connectorconsumption}{in}
\newcommand{\connectorincome}{,}
\newcommand{\balmodifierunb}{}
\newcommand{\balmodifierbal}{, Balanced Panel}

% Variant and hukou names contain underscores, so define via
% \expandafter+\csname (underscores aren't valid in \newcommand names).
\expandafter\newcommand\csname variantnameexp\endcsname{Experience}
\expandafter\newcommand\csname variantnameexp_max\endcsname{Max Experience}
\expandafter\newcommand\csname variantnameexp_sh\endcsname{Experience Share}
\expandafter\newcommand\csname variantnameexp_m_sh\endcsname{Max Experience Share}
\expandafter\newcommand\csname variantnamebirth\endcsname{Urban Birth}
\expandafter\newcommand\csname hukounamerural_first\endcsname{Rural Hukou First}
\expandafter\newcommand\csname hukounameurban_first\endcsname{Urban Hukou First}
\expandafter\newcommand\csname hukounamerural_only\endcsname{Only Rural Hukou}
\expandafter\newcommand\csname hukounameurban_only\endcsname{Only Urban Hukou}

% --- Notes assembly ---
\NewDocumentCommand{\GRCnotesxref}{m m}{%
  \csname datasource#1#2\endcsname{} \refdatasec{} \refIDNcanon. \seclusterind{} \sigstars
}
\newcommand{\GRCnotesIDNcanonical}{%
  \datasourceIDNunb{} \refdatasec.
  The dependent variable is the log of total consumption per capita.
  Urban is an indicator equal to one for individuals who report living in a city or town, as opposed to a village.
  Individuals are assigned to trajectories based on their location history across the survey waves.
  $\Delta_{\text{never}}$ reports the extrapolated returns to migrating to an urban location for individuals who are never observed in an urban location in the data.
  $\bar{\Delta}$ is the sample-weighted average return across switcher trajectories, $\bar{\Delta} = \sum_{\underline{d}} \pi_{\underline{d}} \Delta_{\underline{d}}$, where $\pi_{\underline{d}}$ is trajectory~$\underline{d}$'s share of the estimation sample and $\Delta_{\underline{d}}$ is its trajectory-specific return.
  Columns~(2) to~(5) include time (survey wave) fixed effects, column~(3) adds a female indicator, column~(4) adds age squared, and column~(5) adds education (years of schooling, maximum across periods) and its square.
  \seclusterind{} \sigstars
}
\newcommand{\GRCnotesIncomeShared}{%
  The dependent variable is the log of income per capita.
  Urban is an indicator equal to one for individuals who report living in a city or town, as opposed to a village.
  Individuals are assigned to trajectories based on their location history across the survey waves.
  $\Delta_{\text{never}}$ reports the extrapolated returns to migrating to an urban location for individuals who are never observed in an urban location in the data.
  $\bar{\Delta}$ is the sample-weighted average return across switcher trajectories, $\bar{\Delta} = \sum_{\underline{d}} \pi_{\underline{d}} \Delta_{\underline{d}}$, where $\pi_{\underline{d}}$ is trajectory~$\underline{d}$'s share of the estimation sample and $\Delta_{\underline{d}}$ is its trajectory-specific return.
  Columns~(2) to~(5) include time (survey wave) fixed effects, column~(3) adds a female indicator, column~(4) adds age squared, and column~(5) adds education (years of schooling, maximum across periods) and its square.
  \seclusterind{} \sigstars
}

% --- Main table macros ---
%
% \GRCtable: 5-col main GRC tables (5_GrRC.do, 6_GrRC_NonAg.do).
%   #1=country (IDN/CHN/TZA), #2=depvar (consumption/income),
%   #3=choice (urban/nonag), #4=balance (unb/bal),
%   #5 (optional) = notes-prefix override; defaults to Template B.
%   Use [\GRCnotesIDNcanonical] for the IDN cons/urban/unb canonical
%   table and [\GRCnotesIncomeShared] for income tables.
\NewDocumentCommand{\GRCtable}{m m m m O{}}{%
  \def\GRCthisnotes{#5}%
  \ifx\GRCthisnotes\empty
    \def\GRCthisnotes{\GRCnotesxref{#1}{#4}}%
  \fi
  \begin{table}[htbp]\centering
  \caption{Restricted GRC Estimates of the Returns to \csname treatment#3\endcsname{} on log \csname depvar#2\endcsname{} \csname connector#2\endcsname{} \csname countryname#1\endcsname\csname balmodifier#4\endcsname}%
  \label{tab:GRC_#1_#2_#3_#4}
  \begin{threeparttable}
  \input{tables/GRC_#1_#2_#3_#4}%
  \begin{tablenotes}[flushleft]\footnotesize
    \item \GRCthisnotes
  \end{tablenotes}
  \end{threeparttable}
  \end{table}%
}

% \GRCexptable: 4-col experience and birth family tables (10/11/12/13/14/15).
%   #5=variant (exp/exp_max/exp_sh/exp_m_sh/birth)
\NewDocumentCommand{\GRCexptable}{m m m m m}{%
  \begin{table}[htbp]\centering
  \caption{Restricted GRC Estimates of the Returns to \csname treatment#3\endcsname{} on log \csname depvar#2\endcsname{} \csname connector#2\endcsname{} \csname countryname#1\endcsname\csname balmodifier#4\endcsname, \csname variantname#5\endcsname{} Controls}%
  \label{tab:GRC_#1_#2_#3_#4_#5}
  \begin{threeparttable}
  \input{tables/GRC_#1_#2_#3_#4_#5}%
  \begin{tablenotes}[flushleft]\footnotesize
    \item \csname datasource#1#4\endcsname{} \refdatasec{} \refIDNcanon. \seclusterind{} \sigstars
  \end{tablenotes}
  \end{threeparttable}
  \end{table}%
}

% \GRChukoutable: hukou-subgroup tables (8_GrRC_hukou.do).
%   #5=hukou subgroup (rural_first/urban_first/rural_only/urban_only)
\NewDocumentCommand{\GRChukoutable}{m m m m m}{%
  \begin{table}[htbp]\centering
  \caption{Restricted GRC Estimates of the Returns to Urban Location on log \csname depvar#2\endcsname{} \csname connector#2\endcsname{} \csname countryname#1\endcsname\csname balmodifier#4\endcsname, \csname hukouname#5\endcsname}%
  \label{tab:GRC_#1_hukou_#5_#2_#3_#4}
  \begin{threeparttable}
  \input{tables/GRC_#1_hukou_#5_#2_#3_#4}%
  \begin{tablenotes}[flushleft]\footnotesize
    \item \csname datasource#1#4\endcsname{} \refdatasec{} \refIDNcanon. \seclusterind{} \sigstars
  \end{tablenotes}
  \end{threeparttable}
  \end{table}%
}

% Heterogeneity table macros (Templates D/E) deferred until 16_heterogeneity_tables.do
% canonical IDN heterogeneity prose has been read in detail (Phase 1b.3).

% =====================================================================
% End of GRC table macros block
% =====================================================================